\section{Open Questions}

\begin{enumerate}
  \item \textbf{Hybridization with randomized compilation.} How does the Chebyshev-multiproduct construction compose with qDRIFT or qSWIFT --- does the $O(\log m)$ conditioning survive when the base $U_2$ step is itself replaced by a stochastically-sampled channel, and does the total sample complexity beat qDRIFT alone? \emph{Basis:} the paper is deterministic; qDRIFT achieves $H$-norm-independent scaling via randomization, and the small MPF constant here ($\sim 2$) is now competitive with the qDRIFT variance term. \emph{Next steps:} implement a classical Monte-Carlo simulator composing Chebyshev-MPF ($m=3$) over qDRIFT samples of size $s$; measure diamond-distance-to-target as a function of $(m,s)$ on the same $N=4$ Heisenberg benchmark; compare total gate count to pure qDRIFT and pure MPF-with-Suzuki.
  \item \textbf{Time-dependent Hamiltonians.} Does the well-conditioned coefficient construction generalize to $H(t)$ --- specifically, do the Chebyshev-point $k_j$ still cancel the leading Magnus-expansion terms, or does the time-ordering commutator $[H(t_1),H(t_2)]$ break the Vandermonde structure? \emph{Basis:} paper's proof relies on BCH expansion of a static generator; time-ordered Dyson/Magnus series introduces double-commutators whose cancellation may need different constraints. \emph{Next steps:} derive the time-ordered analog of the Vandermonde system by expanding Magnus to order $2m$; check symbolically whether the Chebyshev-point $k_j$ still zero the coefficients; numerically simulate a driven Heisenberg chain.
  \item \textbf{Hardware-noise sensitivity.} How sensitive is the well-conditioned multiproduct advantage to realistic gate-level noise --- if each $U_2$ base call carries depolarizing rate $p$ per two-qubit gate, at what $m$ does the coherent (algorithmic) error savings get erased by $2m$-fold accumulated incoherent error? \emph{Basis:} paper is noise-free; $m$-fold MPF costs $O(m^2 N)$ gates whose incoherent error scales as $m^2$, competing with the $O((\Delta/r)^{2m})$ coherent gain. \emph{Next steps:} simulate $N=4$ Heisenberg with a per-gate depolarizing channel (stim or qiskit density-matrix); compare Chebyshev-MPF $m=1..6$ vs $U_2$ alone at fixed gate budget; identify $m^*(p)$ for $p \in [10^{-4},10^{-2}]$.
  \item \textbf{Adaptive coefficient sets during runtime.} Can the coefficient set be chosen adaptively at runtime --- picking a different $m$ (and Chebyshev-point set) per step based on a cheap error estimator --- to achieve a tighter overall diamond-norm bound than any fixed-$m$ schedule? \emph{Basis:} paper fixes $m$ offline; adaptive step-size and adaptive order are standard in classical ODE integration (LSODA), with no MPF analog in the literature. \emph{Next steps:} prototype an adaptive controller running $m$ and $m+1$ in parallel with the difference as an embedded error estimator (Fehlberg-style); test on $H(t)$ with sharp features (avoided crossings); compare to fixed-$m$ schedule.
  \item \textbf{Application to concrete chemistry Hamiltonians.} For chemistry Hamiltonians (H-chain, FeMoco active space), does the LCU success probability $1/\|a\|_1^2$ actually deliver the promised $O(t\lambda \log^2(t\lambda/\varepsilon))$ end-to-end gate count, or does the second-quantized fermion overhead swamp the conditioning gain? \emph{Basis:} paper's cost bound is per-MPF-step; end-to-end simulation is often dominated by Jordan-Wigner / Bravyi-Kitaev overhead. \emph{Next steps:} take a small ab-initio Hamiltonian (H$_4$ STO-3G, $\sim 8$ qubits); symbolically construct the Chebyshev-MPF wrapper around a JW Trotter step; count T-gates for target error $10^{-3}$; compare to qubitization; report the $m^*$ at which MPF beats direct qubitization.
\end{enumerate}
